\documentclass[11pt]{article}
\usepackage{amsmath,amssymb,amsthm}
\usepackage[margin=1in]{geometry}
\usepackage{hyperref}
\newcommand{\Q}{\mathbb{Q}}\newcommand{\zetaK}{\zeta_{K}}

\title{A Shared Geometric Object Across the VFD Programme:\\ one object, two roles,
an open bridge}
\author{(VFD programme --- connective note)}
\date{2026}

\begin{document}
\maketitle

\begin{abstract}
This short note records a factual observation and immediately bounds it. The same
geometric object---the $600$-cell and its icosian ring---appears in two different
parts of the VFD programme: \emph{geometrically}, as a substrate in the programme's
cosmological modelling, and \emph{arithmetically}, as the carrier of certain
$L$-functions (companion Paper~I). We state plainly that this is a recurrence of one
\emph{object} in two \emph{roles}; we make \textbf{no claim} that the cosmological
model explains, implies, or is explained by the Riemann zeta function. The bridge
between the geometric and arithmetic roles is \emph{open}, and for the specific case
of the Riemann $\zeta$ it has been tested and found \emph{negative} (the $(O2)$
result, companion Paper~II).
\end{abstract}

\section{The observation}
The $600$-cell / icosian ring is a single, well-defined mathematical object. In the
programme it is used in two distinct ways:
\begin{itemize}
  \item \textbf{geometric role:} as a $4$-dimensional substrate / vertex
  configuration in cosmological modelling;
  \item \textbf{arithmetic role:} as the order whose theta series realizes
  $\zetaK(s)\zetaK(s-1)$, with sub-objects realizing $\zeta(s)\zeta(s-1)$ and
  $L(\chi_5)$ (Paper~I).
\end{itemize}
That one object recurs in both is true and worth recording. It is the natural reason
the same vocabulary (spheres, polytopes, icosahedral symmetry) appears in both
settings.

\section{What is \emph{not} claimed}
We are explicit, because the tempting overstatement is exactly the one to avoid.
\begin{enumerate}
  \item We do \textbf{not} claim that the cosmological model is, encodes, or
  explains the Riemann $\zeta$ or its zeros.
  \item We do \textbf{not} claim a proven bridge between the geometric and arithmetic
  roles of the object. That bridge is \emph{open}.
  \item For the specific, strongest form of such a bridge---``the substrate realizes
  the Riemann $\zeta$''---the engine of Paper~II returns a \emph{certified negative}
  on five independent grounds: the substrate realizes $\zetaK=\zeta\cdot L(\chi_5)$,
  which contains $\zeta$ only as a factor.
\end{enumerate}

\section{Two cautionary anchors}
Two elementary facts guard against reading more into the recurrence than is there.
\begin{itemize}
  \item \textbf{Shared silhouette is not shared content.} Diagrams in cosmology and
  in analytic number theory can share an outline (e.g.\ a flaring/funnel shape)
  while denoting unrelated objects; the visual rhyme carries no implication.
  \item \textbf{Universality is generic.} The Riemann zeros share their fine
  statistics (GUE) with a vast range of unrelated systems; sharing a universality
  class is therefore \emph{weak} evidence of any specific identity.
\end{itemize}

\section{Statement of record}
\emph{One object recurs in two roles; the cross-role bridge is open; the specific
$\zeta$-bridge is certified negative; no claim links the cosmological model to the
Riemann Hypothesis.} This note exists so that the recurrence can be cited accurately
without being inflated.

\end{document}
